\documentclass[instructions.tex]{subfiles}
\begin{document}
 
\chapter{Du scandale et du mauvais exemple.}
 
\primo IL n’y a rien que nous devions plus craindre que de scandaliser, c’est-à-dire, de faire pécher une âme : Malheur à l’homme par qui le scandale arrive ! Il ne faut qu’une parole, qu’un regard, qu’une parure, qu’une chanson, qu’un conseil, qu’une flatterie, qu’un rapport, pour faire pécher les autres : Une étincelle de feu peut causer un grand embrasement.
 
Notre charité doit nous engager à ménager même la délicatesse des esprits foibles, et à nous priver de certaines choses indifférentes ou permises, de certaines visites, de certaines manières de vivre, lorsque quelqu’un s’en scandalise ; si je savais, dit saint Paul, que mon frère se scandalisât de me voir manger de la chair, je n’en mangerais jamais, de crainte de scandaliser mon frère\footnote{Si esca scandalizat fratrem meum, non manducabo carnem in æternum, ne fratrem meum scandalizem. 1. Cor. 8. 13.} Il vaudrait mieux, dit Jésus-Christ, être précipité dans la mer, que de scandaliser même un enfant.
 
\secundo Si c’est un mal de scandaliser dans des choses indifférentes, que sera-ce dans des choses criminelles ? O filles vaines ! de combien de péchés de pensées, de regards et de désirs criminels n’êtes-vous pas la cause par vos complaisances, vos manières libres et vos airs mondains ! Libertins, combien de péchés n’ont pas fait commettre vos plaisanteries sur la religion et sur la pudeur, vos sollicitations au crime, à la débauche ! Bouches envenimées, de combien de divisions et de querelles avez-vous été la cause par vos médisances et vos rapports flatteurs ! En réparerez-vous les suites funestes ?
 
O vous qui, par votre autorité, devez empêcher les désordres de votre maison ; ou qui, par votre emploi, devez réprimer les désordres publics, de combien de crimes n’êtes-vous pas responsables, si vous négligez vos devoirs ! Et vous, riches sans humanité, comment répondrez-vous à Dieu du scandale des pauvres qui vous chargent de malédictions, murmurent contre le ciel, voyant qu’ils manquent de tout, tandis que vous ne manquez de rien ? Pères et mères, qui ne vivez pas saintement, comprenez-vous quel tort vous faites à vos enfans par vos exemples ? Il serait à souhaiter que jamais ils ne fussent nés de vous ; vous ne leur avez donné la vie du corps que pour donner à leur âme la mort éternelle.
 
Les mauvais exemples ont perdu plus d’âmes que les bons n’en ont sauvé. Si nous entendions les plaintes des réprouvés, il n’y en a pas un qui ne nous dit avec des cris lamentables : C’est une telle, c’est un tel qui m’a perdu par les mauvais conseils qu’il m’a donnés, par les mauvais livres qu’il m’a prêtés, par les fréquentations, les débauches et les compagnies où il m’a engagé.
 
Dieu nous commande d’aimer même ceux qui nous font du tort ; pourquoi faisons-nous périr des âmes qui ne nous ont jamais fait de mal ? A quel jugement devez-vous vous attendre, quand vous paraîtrez au tribunal de Jésus-Christ ? Comment obtiendrez-vous sa miséricorde, après lui avoir enlevé des âmes qui lui ont coûté son sang, et pour lesquelles il a donné sa vie ?
 
Celui qui, par sa faute, donne aux autres, par conseils, par écrit, ou par sa conduite, occasion prochaine de se pervertir, est responsable, dit Salvien, de tous les désordres dont il a été la cause, et souffrira autant d’enfers qu’il aura fait de malheureux\footnote{Pro tantis reus, quantos secum traxerit in reatum. Salv.} : c’est-à-dire que vous rendrez compte à Dieu, âme pour âme, de celles que vous aurez perdues, et de celles dont vous étiez chargé ; que ces âmes, perdues par votre faute, demanderont à Dieu vengeance contre vous, et seront dans l’éternité vos plus cruels bourreaux. Vous éviterez ce malheur, si, par vos exemples et vos prières, vous tâchez de sauver autant d’âmes que vous en avez perverti, et d’empêcher, s’il se peut, autant de péchés que vous en avez fait commettre.
 
Lorsque votre conduite est irréprochable, que vous vaquez à votre devoir, si quelqu’un s’en scandalise, ce n’est plus votre faute, vous n’en répondrez pas ; ce n’est pas un scandale que vous lui donnez, mais un scandale qu’il prend par sa malice. Tels étaient les pharisiens, qui se scandalisaient de la conduite de Jésus-Christ ; tels sont aujourd’hui les voluptueux et les impies, qui se raillent et se scandalisent de la vertu et du zèle des pasteurs.
 
\section{Résolutions}
 
\primo Je craindrai de donner jusqu’au moindre scandale. 

\secundo Je réparerai ceux que j’ai pu donner jusqu’ici, par une conduite édifiante en tout. 

\tertio Et je prierai tous les jours pour les personnes que j’ai eu le malheur de scandaliser. 

\prayer{Mon Dieu, pardonnez-moi mes misérables scandales ; ayez pitié des personnes que j’ai portées au mal ; daignez leur pardonner les fautes dont j’ai été la cause, et ramenez-les au bien.}
 
\end{document}
